{\setlength{\parskip}{8mm}
{\huge{\textbf{Resumen}}}

La estructura de este documento es orientativa. Las \textbf{partes obligatorias (motivación y objetivos, competencias, aportaciones, alineamiento ODS, desarrollo, conclusiones y bibliografía)} deben estar recogidas, aunque pueden reordenarse o redistribuirse en base a las características del trabajo.

    {EL presente Trabajo de Fin de Título(TFT) desarrolla un sistema empotrado basado en bus CAN que trabaja en la capa de aplicación de OBDII permitiendo la realización de diagnósis de vehículos. El sistema se ha implementado utilizando como vehículo de pruebas un SEAT Ibiza 6j con motor 1.4 BXW(\textit{sustituir codigo motor por ECU}).}

\newpage

{\huge{\textbf{Abstract}}}


    {\large{Ahora en inglés: 
    This Final Degree Project (TFT) develops an embedded system based on the CAN bus that operates at the OBD-II application layer, enabling vehicle diagnostics. The system has been implemented using a SEAT Ibiza 6J with a 1.4 BXW engine (\textit{replace engine code with ECU designation}) as the test vehicle.
 }

}}